% \section{Method}
% \label{sec:method}

% \subsection{Temporal Grounding}
% \label{sec:method:grounding}

% The temporal grounding stage transforms a natural-language temporal question into its structured reasoning representation, establishing a grounded basis for subsequent retrieval and reasoning. 
% It comprises three key phases: temporal fact grounding, question analysis, and temporal question decomposition. 
% An overview of the entire temporal grounding workflow is illustrated in Figure~\ref{fig:initial}.

% \myparagraph{Knowledge fact grounding}
% Given an input question $Q$, \mot first identifies a relevant temporal subgraph $\mathcal{G}_Q$ within the underlying TKG.
% This subgraph captures all entities, relations, and time-stamped triples that are semantically or temporally related to $Q$ within a bounded $D_{\max}$-hop neighborhood, 
% serving as the factual grounding for subsequent reasoning.

% \myparagraphunderline{Topic entity recognition}
% To locate question-relevant entities, \mot employs LLMs to extract potential entity mentions and associated timestamps from $Q$. 
% Following extraction, a Dense Retrieval Model (DRM) aligns these mentions with KG entities using embedding-based similarity matching. 
% Specifically, both question keywords and KG entities are encoded into dense embeddings, and an entity index is constructed using FAISS. 
% A cosine similarity matrix is then computed between the two embedding spaces, and the highest-scoring entities are selected to form the topic entity set $Topics$. 
% These topics act as initial anchors for building the question-specific subgraph and initializing temporal retrieval.

% \myparagraphunderline{Subgraph construction}
% After identifying the topic entities, \mot constructs a $D_{\max}$-hop evidence subgraph $\mathcal{G}_Q$ that includes all triples connected to each topic entity within the defined temporal window. 
% This subgraph provides a localized yet sufficiently rich factual context for reasoning. 
% To mitigate redundancy and computational overhead, we apply graph reduction and relation clustering techniques following~\cite{pogtan2025paths}, ensuring $\mathcal{G}_Q$ remains compact and semantically coherent. 
% The resulting subgraph forms the temporal evidence foundation for all subsequent modules.


% \myparagraph{Question analysis}
% Once the temporal subgraph $\mathcal{G}_Q$ is obtained, the question is analyzed to infer its underlying temporal intent. 
% Unlike conventional KGQA approaches that treat all questions uniformly, temporal questions often involve diverse forms and complex temporal constraints. 
% Therefore, \mot first performs temporal type classification to identify the temporal operator implied in $Q$ before proceeding to question decomposition.

% \myparagraphunderline{Temporal type classification.}
% Temporal type classification determines the intrinsic temporal relations expressed in a question. 
% Following Allen’s interval algebra~\cite{allen1984towards}, we consider thirteen fundamental temporal relations (\eg \textsc{before}, \textsc{meets}, \textsc{overlaps}, \textsc{during}, \textsc{starts}, \textsc{finishes}, \textsc{equals}), 
% as well as extended temporal operators such as set relations, duration comparisons, and ordering operations (\eg \textsc{first}, \textsc{last}, \textsc{n-th})~\cite{TimelineKGQA}.  
% Existing methods typically rely on static prompt templates or manually curated exemplars, which limits their adaptability to diverse temporal patterns.  
% In contrast, \mot employs a continually updated experience pool $\mathcal{E}_{pool}$ that stores successful reasoning trajectories and LLM operations. 
% This memory structure enables dynamic retrieval of contextually relevant exemplars to enhance in-context classification.
% Formally, for a given question $Q$, type-specific exemplars are retrieved as:
% \[
% \mathcal{E}_{\text{type}} = \textsf{GetTypeExp}(\mathcal{E}_{pool}, Q, W_{\text{exp}}),
% \]
% where $W_{\text{exp}}$ denotes the exemplar retrieval limit. 
% These exemplars are incorporated into the classification prompt as:
% \[
% \textsf{Type}(Q) = \textsf{Prompt}_{\text{TypeSelect}}(Q, \mathcal{E}_{\text{type}}),
% \]
% allowing the LLM to infer the most appropriate temporal operator. 
% This adaptive, exemplar-guided strategy improves type prediction accuracy and ensures generalization to unseen question forms.




\section{Method}
\label{sec:method}

% \noindent\textbf{Overview.}
% \mot\ implements \emph{temporal KG–grounded LLM reasoning} by (i) grounding a natural-language question on a question-specific temporal subgraph, (ii) recursively decomposing the question into executable sub-queries with experience-guided control, (iii) executing experience-informed toolkits to retrieve and prune temporally valid evidence paths, and (iv) writing verified traces back to a continually evolving experience memory. 
% Unlike prior methods that rely on a single retriever or fixed prompts, \mot\ couples a temporal-first pruning policy with an experience layer that adapts prompts, toolkit choices, and reuse decisions over time. 
% The overall framework (Figure~\ref{fig:overall}) comprises four components:

% \begin{itemize}[leftmargin=*]
%   \item \textbf{Temporal Grounding.} 
%   From the input question, \mot\ links topic entities and constructs a $D_{\max}$-hop temporal evidence subgraph $\mathcal{G}_Q$. 
%   It then performs temporal type classification (via experience-augmented prompts) and generates a hierarchical question tree with explicit indicators $\Ind_i=\langle x?,R,y?,\tau;\mathcal{C}_{time}\rangle$ that normalize operators, bounds, and granularity for downstream execution (Section~\ref{sec:method:grounding}).

%   \item \textbf{Recursive Reasoning (Controller).} 
%   Given the question tree, \mot\ follows a memory-first policy: for each node $q_i$, it reuses verified trajectories when sufficient; otherwise, it selects one or more temporal toolkits using experience-guided prompts. 
%   The controller coordinates sufficiency testing, lightweight debate–vote aggregation, and—if needed—refinement or further decomposition under depth/branch limits, progressively assembling a coherent reasoning chain (Section~\ref{sec:recursion}).

%   \item \textbf{Temporal Facts Retrieval \& Pruning (Execution).} 
%   The selected toolkits trigger a hybrid retriever that combines time-monotone graph exploration with dense embedding lookup. 
%   Candidates are first filtered by strict temporal validity (constraint satisfaction and monotonicity), then re-ranked by semantic compatibility and temporal proximity; a final LLM-aware selector retains Top-$W_{\max}$ evidence paths for verification (Section~\ref{sec:retrieval}).

%   \item \textbf{Experience Memory.} 
%   Verified sub-answers, paths, indicators, and toolkit decisions are encoded with dual embeddings (question and indicator) and stored together with type labels and hit statistics. 
%   At query time, retrieval is scoped by type and ranked by a hybrid score (similarity and recency/hit-count), while write-back continuously refreshes the exemplar pool and short-term buffer—closing the retrieval–reasoning–learning loop (Section~\ref{sec:memory}).
% \end{itemize}
\noindent\textbf{Overview.}
\mot\ implements ``TKG-grounded LLM reasoning'' by grounding each question in temporal facts, recursively decomposing it into executable sub-queries, retrieving and pruning temporally valid evidence, and continually refining its experience memory. 
Unlike prior TKGQA or RAG systems that depend on static retrievers or fixed prompts, \mot\ integrates temporal alignment, hierarchical control, and self-evolving memory into a unified reasoning framework. 
The overall architecture of \mot is detailed in Figure~\ref{fig:overall}, consisting of four key components:

\begin{itemize}[leftmargin=*]
  \item \textbf{Temporal Grounding.} 
  The model begins by linking topic entities and constructing a $D_{\max}$-hop temporal subgraph $\mathcal{G}_Q$ from the TKG. 
  It then classifies the temporal type of the question using exemplars retrieved from memory, producing structured temporal constraints for downstream reasoning (Section~\ref{sec:method:grounding}).

  \item \textbf{Tree of Time Reasoning.} 
  Given the grounded question and its temporal type, \mot\ constructs a hierarchical decomposition tree and executes sub-questions in a top-down manner. 
  For each node, it dynamically decides whether to recall prior experiences, invoke temporal toolkits, or refine unresolved branches, ensuring temporal consistency and interpretability (Section~\ref{sec:recursion}).

  \item \textbf{Temporal Evidence Retrieval and Pruning.} 
  Guided by toolkit configurations, \mot performs hybrid retrieval, combining time-monotone, graph exploration, and embedding-based search. 
  Retrieved candidates are filtered by temporal constraints, re-ranked by semantic and temporal proximity, and finally verified through an LLM-aware selection step (Section~\ref{sec:retrieval}).

  \item \textbf{Experience Memory.} 
  Verified reasoning traces, toolkit selections, and sub-question embeddings are stored in a dynamic memory pool. 
  At inference, similar experiences are retrieved by type and relevance; after reasoning, new traces are written back to continuously refine retrieval and decision-making across future questions (Section~\ref{sec:memory}).
\end{itemize}
\vspace{-4mm}
\subsection{Temporal Grounding}
\label{sec:method:grounding}

The temporal grounding stage transforms a natural-language temporal question into a structured reasoning representation, establishing a factual and temporal foundation for subsequent retrieval and inference.  
It consists of two core phases: knowledge fact grounding and question analysis.  The pseudo-code of the temporal grounding is detailed in Algorithm \ref{algorithm:analysis} of Appendix \ref{appendix:alg}.
% An overview of the entire grounding workflow is shown in Figure~\ref{fig:initial}.

\myparagraph{Knowledge fact grounding}
Given an input question $Q$, \mot first identifies the relevant temporal subgraph $\mathcal{G}_Q$ within the underlying TKG.  
This subgraph captures entities, relations, and time-stamped triples that are semantically or temporally associated with $Q$ within a bounded $D_{\max}$-hop neighborhood, 
serving as the factual evidence base for downstream reasoning.

\myparagraphunderline{Topic entity recognition}
To locate question-relevant entities, \mot employs LLMs to extract potential entity mentions and associated timestamps from $Q$.  
Following extraction, a Dense Retrieval Model (DRM) aligns these mentions with KG entities via embedding-based similarity matching.  
Specifically, both question keywords and KG entities are encoded into dense embeddings, and an entity index is constructed using FAISS \cite{douze2024faiss}.  
Cosine similarity is then computed between the two embedding spaces, and the top-ranked entities are selected to form the topic entity set $Topics$.  
These topic entities act as the initial anchors for constructing a localized subgraph.

\myparagraphunderline{Subgraph construction}
Once the topic entities are identified, \mot constructs a $D_{\max}$-hop subgraph $\mathcal{G}_Q$ that aggregates all triples connected to each topic entity within the defined temporal window.  
This subgraph provides a compact yet semantically rich neighborhood for temporal reasoning.  
To reduce redundancy and computational overhead, we apply graph reduction and relation clustering techniques following~\cite{pogtan2025paths}, ensuring $\mathcal{G}_Q$ remains concise while preserving the most relevant temporal relations.  
% The resulting evidence subgraph serves as the contextual grounding for all subsequent modules.

\myparagraph{Question analysis}
After constructing $\mathcal{G}_Q$, the next step is to analyze the temporal intent of the question.  
Unlike conventional TKGQA approaches that treat all questions uniformly, or use the label from the dataset, leveraging the flexibility reduces. Temporal questions exhibit diverse syntactic forms and intricate time constraints.  
To capture these nuances, \mot first performs {temporal type classification}, identifying the temporal operator that governs the question’s reasoning logic.

\myparagraphunderline{Temporal type classification}
Temporal type classification determines the intrinsic temporal relations implied by a question.  
Following Allen’s interval algebra~\cite{allen1984towards}, we consider thirteen fundamental temporal relations,  
as well as extended operators such as set relations, duration comparisons, and ordering mechanisms~\cite{TimelineKGQA}, as introduced in Section \ref{sec:preliminaries}.  
Existing methods often rely on static prompts or manually designed exemplars, limiting their adaptability to complex or unseen temporal structures.  
In contrast, \mot leverages a continuously updated experience pool $\mathcal{E}_{pool}$ that stores successful reasoning trajectories, toolkit configurations, and classification results.  
This memory-based design enables dynamic retrieval of contextually aligned exemplars, improving temporal reasoning consistency across question types.  
Formally, for a given question $Q$, type-specific exemplars are retrieved as:
$
\mathcal{E}_{\text{type}} = \textsf{GetTypeExp}(\mathcal{E}_{pool}, Q, W_{\text{exp}}),
$
where $W_{\text{exp}}$ denotes the exemplar retrieval limit.  
These exemplars are then incorporated into the LLM classification prompt:
$
\textsf{Type}(Q) = \textsf{Prompt}_{\text{TypeSelect}}(Q, \mathcal{E}_{\text{type}}),
$
allowing the model to identify the most appropriate temporal operator for $Q$.  
% The prompt template is presented in Appendix \ref{appendix:prompt}.
This exemplar-guided strategy enables the \mot to adapt to new question patterns, enhances type prediction accuracy, and lays the foundation for the subsequent decomposition stage.
